\chapter{Công nghệ mạng trong Internet vạn vật}
\label{ch:network}
\begin{keybox}
\textbf{Mục tiêu học tập.} Sau chương này, người học giải thích được H2M/M2M/M2H/H2H; phân biệt tên, định danh, địa chỉ IPv4/IPv6; chọn được công nghệ kết nối phù hợp theo phạm vi--băng thông--năng lượng; trình bày bảo mật và định tuyến RPL; phân tích HTTP, CoAP, MQTT, WebSocket, AMQP, XMPP và vận dụng để đề xuất kiến trúc IoT.
\end{keybox}
\noindent\textit{Quy ước nguồn:} ``Theo slide'' chỉ nội dung xuất hiện trong giáo trình giảng dạy; ``Mở rộng kỹ thuật'' là giải thích và chuẩn hóa bổ sung nhằm hỗ trợ hiểu bản chất. Các sơ đồ dưới đây được vẽ lại độc lập.
\section{Giao tiếp trong IoT (Communications in IoT)}
\subsection{Khái niệm và vai trò}
Internet of Things (IoT) là tập hợp các thực thể vật lý có cảm biến, bộ xử lý, phần mềm và khả năng kết nối để thu nhận, truyền, nhận dữ liệu và tác động lên môi trường. Một giải pháp IoT thường gồm \,\emph{thiết bị đầu cuối} (end device), \emph{liên kết truyền thông}, \emph{gateway/broker}, \emph{dịch vụ xử lý} và \emph{giao diện người dùng}. Thiết bị có thể hoạt động tự động không cần tương tác con người ở mỗi lần truyền dữ liệu. Không phải mọi giao tiếp cảm biến đều trực tiếp đi ra Internet: thiết bị dùng BLE, Zigbee hoặc RF có thể gửi qua gateway.\srcslide{2--3}
\begin{figure}[htbp]\centering
\begin{tikzpicture}[>=Latex,node distance=7mm and 7mm,font=\footnotesize]
\node[draw,rounded corners,fill=pale,minimum width=25mm,minimum height=13mm] (s) {Cảm biến};
\node[draw,rounded corners,fill=pale,right=of s,minimum width=25mm,minimum height=13mm] (m) {Vi điều khiển};
\node[draw,rounded corners,fill=pale,right=of m,minimum width=25mm,minimum height=13mm] (g) {Gateway};
\node[draw,rounded corners,fill=blue!10,right=of g,minimum width=25mm,minimum height=13mm] (c) {Cloud/Broker};
\draw[->,thick,blue] (s)--(m); \draw[->,thick,blue] (m)--node[above]{BLE/RF}(g); \draw[<->,thick,blue] (g)--node[above]{IP}(c);
\node[draw,rounded corners,fill=accent!12,below=10mm of c,minimum width=25mm,minimum height=12mm] (u) {Ứng dụng};
\draw[<->,thick] (c)--(u);
\node[draw,rounded corners,fill=accent!12,below=10mm of m,minimum width=25mm,minimum height=12mm] (a) {Cơ cấu chấp hành};
\draw[->,thick] (m)--(a);
\end{tikzpicture}
\caption{Sơ đồ tham khảo một luồng dữ liệu IoT từ cảm biến đến ứng dụng và điều khiển ngược.}
\end{figure}
\subsection{Bốn kiểu giao tiếp}
\begin{description}[style=nextline,leftmargin=6mm]
\item[Human-to-Machine (H2M) — Người đến máy.] Người cung cấp lệnh qua tiếng nói, văn bản, hình ảnh, nút bấm; thiết bị phân tích và hành động. Ví dụ: điều khiển loa thông minh bằng giọng nói, xác thực vân tay, điểm danh khuôn mặt. Đầu vào thuộc về con người, hệ thống là bên tiếp nhận.
\item[Machine-to-Machine (M2M) — Máy đến máy.] Hai hoặc nhiều máy tự trao đổi dữ liệu/lệnh không cần người thao tác trong từng giao dịch; có thể truyền có dây hoặc vô tuyến. Ví dụ: máy bán hàng báo hết hàng cho hệ thống hậu cần, đồng hồ điện báo chỉ số, máy giặt gửi trạng thái. IoT kế thừa M2M nhưng có thể mở rộng tới hệ sinh thái dịch vụ và cloud.
\item[Machine-to-Human (M2H) — Máy đến người.] Thiết bị chủ động thông báo bằng còi, đèn, tin nhắn, âm thanh, hình ảnh. Ví dụ: báo cháy, đèn giao thông và thiết bị theo dõi sức khỏe cảnh báo bất thường.
\item[Human-to-Human (H2H) — Người đến người.] Người trao đổi thông tin để phối hợp điều hành, xử lý bất thường hoặc quyết định. Trong hệ thống IoT, báo cáo từ máy có thể dẫn đến sự phối hợp giữa kỹ thuật viên, nhân viên giám sát và người vận hành.
\end{description}
\srcslide{3--7}
\begin{table}[htbp]\centering\small
\caption{Đối chiếu các dạng giao tiếp theo hướng phát động}
\begin{tabularx}{\textwidth}{@{}p{15mm}X X@{}}\toprule
\textbf{Loại} & \textbf{Ai khởi tạo thông tin?} & \textbf{Ví dụ}\tabularnewline\midrule
H2M & Con người phát lệnh cho thiết bị & Ra lệnh bật đèn bằng giọng nói\\
M2M & Thiết bị tự trao đổi với thiết bị/dịch vụ & Đồng hồ điện truyền số đọc\\
M2H & Thiết bị chủ động cảnh báo cho người & Cảnh báo khói từ cảm biến\\
H2H & Con người trao đổi để phối hợp & Kỹ thuật viên xử lý cảnh báo\\\bottomrule
\end{tabularx}
\end{table}
\section{Định danh, đặt tên và địa chỉ (Identification and Addressing)}
\subsection{Phân biệt định danh, tên và địa chỉ}
\term{Định danh} (identifier) nhằm nhận biết một đối tượng; \term{tên} (name) biểu diễn đối tượng trong không gian tên; \term{địa chỉ} (address) giúp tìm vị trí hoặc định tuyến liên lạc đến giao diện mạng. Một thiết bị có thể có serial, MAC, địa chỉ IP và URI cùng lúc. Địa chỉ IP không tự thay thế cơ chế xác thực thiết bị; biết định danh cũng không chứng minh một nút đáng tin cậy.\srcslide{9--10}
\begin{itemize}
\item \term{URI (Uniform Resource Identifier):} định danh tài nguyên; ví dụ \texttt{https://example.com/temperature/actual} biểu diễn tài nguyên nhiệt độ. URL là một loại URI chứa cách truy cập.
\item \term{EPC (Electronic Product Code):} mã nhận diện sản phẩm, hay dùng với nhãn RFID, ví dụ \texttt{urn:epc:id:sgtin:0614141.012345.62852}.
\item \term{uCode:} hệ thống mã nhận diện đối tượng, xuất hiện trong slide như một phương án đặt tên khác.
\item \term{MAC address:} định danh giao diện trong nhiều mạng lớp liên kết; \term{IP address:} địa chỉ lớp mạng, có thể thay đổi theo cấu hình mạng.
\end{itemize}
\subsection{IPv4: địa chỉ 32 bit}
IPv4 dùng \(32\) bit, chia thành bốn octet (8 bit) dạng thập phân có dấu chấm: \texttt{192.168.1.10}. Số tổ hợp lý thuyết \(2^{32}=4\,294\,967\,296\); không phải tất cả đều là địa chỉ công cộng có thể cấp cho host. CIDR \texttt{/24} chỉ 24 bit mạng, còn 8 bit host. Với mạng \texttt{192.168.1.0/24}, địa chỉ mạng là \texttt{.0} và broadcast là \texttt{.255}; cấu hình phổ biến dùng \texttt{.1} tới \texttt{.254} cho host.\srcslide{11}
\begin{figure}[htbp]\centering
\begin{tikzpicture}[font=\footnotesize]
\foreach \x/\v/\b in {0/192/11000000,1/168/10101000,2/1/00000001,3/10/00001010}{
\node[draw,minimum width=31mm,minimum height=9mm,fill=blue!10] at (\x*3.15,0) {\texttt{\v}};
\node[draw,minimum width=31mm,minimum height=9mm] at (\x*3.15,-1.15) {\texttt{\b}};
\draw[->] (\x*3.15,-.45)--(\x*3.15,-.7);
}
\node[anchor=east] at (-1.65,0) {Thập phân}; \node[anchor=east] at (-1.65,-1.15) {Nhị phân};
\end{tikzpicture}\caption{Địa chỉ IPv4 gồm bốn trường 8 bit (hình tự dựng).}
\end{figure}
\subsection{IPv6: địa chỉ 128 bit}
IPv6 dùng \(128\) bit, biểu diễn bằng tám nhóm 16 bit dạng hệ 16, ngăn bởi dấu hai chấm. Ví dụ \texttt{2001:0db8:00a0:fe01:0000:0000:0000:0001} có thể viết ngắn \texttt{2001:db8:a0:fe01::1}. Ký hiệu \texttt{::} thay một chuỗi nhóm 0 liên tiếp, chỉ dùng tối đa một lần trong một địa chỉ. Không gian địa chỉ lý thuyết bằng \(2^{128}\approx3.4\times10^{38}\). Tự cấu hình SLAAC và định tuyến phân cấp làm IPv6 hữu ích với IoT quy mô lớn.\srcslide{12--16}
\begin{table}[htbp]\centering\small\caption{IPv4 so với IPv6 theo các đặc trưng cơ bản}
\begin{tabularx}{\textwidth}{@{}p{33mm}X X@{}}\toprule
\textbf{Tiêu chí}&\textbf{IPv4}&\textbf{IPv6}\\\midrule
Độ dài & 32 bit & 128 bit\\
Biểu diễn & 4 số thập phân & 8 nhóm hệ 16\\
Tự cấu hình & DHCPv4 hoặc thủ công & SLAAC, DHCPv6 hoặc thủ công\\
Quảng bá & Có broadcast & Không có broadcast IP, dùng multicast\\
Mạng IoT hạn chế & Có thể dùng qua gateway & Có thể dùng với 6LoWPAN để nén header\\\bottomrule
\end{tabularx}
\end{table}
\begin{note}
Slide nêu \emph{IPv6 mặc định an toàn hơn IPv4} và \emph{gửi gói lớn giúp tiết kiệm băng thông}. Đây không phải tính chất bảo mật hay tốc độ được đảm bảo tự động. IPv6 có hỗ trợ IPsec trong bộ giao thức, nhưng triển khai bảo mật phụ thuộc cấu hình thực tế; hiệu suất còn tùy MTU, overhead, đường truyền và cơ chế ứng dụng. Cần phân biệt ưu điểm của không gian địa chỉ với tính năng bảo mật của giao thức đi kèm.
\end{note}
\subsection{Mô hình phân tầng giao thức}
Kết nối vật lý (radio, cáp, quang) không đồng nghĩa với giao thức ứng dụng (HTTP, MQTT). Một gói MQTT điển hình có thể đặt trên TCP/IP qua Wi-Fi/Ethernet/cellular, còn CoAP thường hoạt động trên UDP/IP. Gateway chuyển tiếp giữa một mạng phi IP ở thiết bị với mạng IP ở phía server.\srcslide{8, 17, 55--59}
\begin{figure}[htbp]\centering
\begin{tikzpicture}[font=\small]
\foreach \y/\a/\b in {3/{Ứng dụng}/{HTTP, MQTT, CoAP, WebSocket},2/{Truyền tải}/{TCP, UDP; TLS/DTLS},1/{Mạng}/{IPv4, IPv6, RPL},0/{Liên kết/Vật lý}/{Wi-Fi, Ethernet, BLE, Zigbee, LoRa}} {
\node[draw,fill=blue!10,minimum width=37mm,minimum height=10mm] at (0,\y*.92) {\a};
\node[draw,fill=white,minimum width=91mm,minimum height=10mm] at (6.8,\y*.92) {\b};
}
\end{tikzpicture}\caption{Góc nhìn phân lớp để tránh nhầm công nghệ kết nối và giao thức dữ liệu.}
\end{figure}
\section{Các công nghệ truyền thông trong IoT}
\subsection{Tiêu chí lựa chọn}
Không có chuẩn vô tuyến nào tối ưu đồng thời về \textbf{phạm vi}, \textbf{thông lượng}, \textbf{mức tiêu thụ điện}, \textbf{độ trễ}, \textbf{chi phí}, \textbf{khả năng xuyên vật cản}, \textbf{độ tin cậy} và \textbf{hạ tầng}. Cần phân biệt \emph{tốc độ PHY} với \emph{thông lượng ứng dụng}, khoảng cách quảng cáo với khoảng cách thực địa. Các khoảng cách và bit rate trong slide 41 chỉ là minh họa, không phải cam kết cho mọi cấu hình.\srcslide{8, 18--41}
\subsection{Truyền thông trên đường dây điện: PLC và X10}
\term{PLC (Power-Line Communication)} truyền tín hiệu dữ liệu trên hệ thống dây cấp điện. Vì thiết bị cố định đã được nối lưới điện, PLC giảm nhu cầu đi thêm dây mạng và có thể hữu ích tại công trình nhiều nhiễu RF hoặc hệ thống điện thông minh. Tuy nhiên chất lượng phụ thuộc tải, lọc nhiễu, topology, bộ biến áp và chuẩn PLC cụ thể. \term{X10} là công nghệ điều khiển nhà thông minh đời đầu dùng xung tín hiệu gần điểm qua không trên đường điện; có tốc độ rất thấp và ít chức năng bảo mật/ACK.\srcslide{18--20}
\begin{table}[htbp]\centering\small\caption{Thông số X10 được nêu trong bài giảng (slide 20)}
\begin{tabularx}{\textwidth}{@{}p{45mm}X@{}}\toprule
Đặc điểm & Thông tin trong slide\\\midrule
Môi trường & Powerline\\
Điều chế & OOK, sóng mang 120 kHz\\
Tốc độ & 120 bps tức thời; 60 bps liên tục; 24 bps payload (cách diễn đạt của slide)\\
Địa chỉ & 16 house code (A--P) \(\times\) 16 unit code (1--16) = 256 tổ hợp\\
Lệnh & ON, OFF, DIM, BRIGHT, ALL UNITS OFF, v.v.\\
Bảo mật/ACK & Slide ghi không hỗ trợ\\\bottomrule
\end{tabularx}
\end{table}
\begin{note}Bảng nguồn có câu ``100 bits (8.5 bytes)'' về độ dài thông điệp. Hai số này mâu thuẫn vì \(100/8=12.5\) byte. Tài liệu không sử dụng chúng như một phép quy đổi chính xác, cần đối chiếu đặc tả định dạng frame cụ thể.\end{note}
\subsection{NFC và UWB}
\term{NFC (Near-Field Communication)} cho phép trao đổi ở cự ly rất ngắn, điển hình vài cm, dùng sóng mang \(13.56\,\mathrm{MHz}\). Theo slide có các tốc độ 106, 212, 424 kbps; các chế độ reader/writer, card emulation và peer-to-peer. Ứng dụng: thanh toán, vé/khóa điện tử, nhận dạng thiết bị, ghép nối ban đầu. Cần lưu ý ``NFC không phải RFID'' là cách diễn đạt giản lược: NFC kế thừa các họ chuẩn RFID tần số cao nhưng mở rộng phương thức giao tiếp chủ động.
\term{UWB (Ultra-Wideband)} dùng tín hiệu trải trên băng tần rất rộng, có thể ứng dụng xác định khoảng cách/thời gian lan truyền và định vị tài sản trong nhà. Đặc trưng là độ phân giải thời gian cao; chất lượng đo vẫn phụ thuộc cấu hình kênh, vật cản, NLOS và đồng bộ. UWB không phải lựa chọn mặc định cho các cảm biến định kỳ chỉ gửi dữ liệu nhỏ.\srcslide{21--23}
\subsection{Wi-Fi, Wi-Fi Direct và Wi-Fi HaLow}
\term{Wi-Fi} (IEEE 802.11) thích hợp thiết bị dùng nguồn ổn định, yêu cầu tải IP trực tiếp hoặc dữ liệu lớn; phổ biến ở 2.4/5 GHz tùy chuẩn, có phiên bản băng tần khác. Wi-Fi Direct tạo liên kết trực tiếp giữa thiết bị mà không bắt buộc dùng access point truyền thống. Đánh đổi phổ biến của Wi-Fi là dòng tiêu thụ khi phát và thời gian kết nối lại. \term{Wi-Fi HaLow} (802.11ah) dùng phổ dưới 1 GHz, nhắm mục tiêu phạm vi xa và tiết kiệm năng lượng hơn trong một số cấu hình cảm biến; cần thiết bị/hạ tầng tương thích, không nên mặc định mọi Wi-Fi là HaLow.\srcslide{24--25}
\subsection{Bluetooth Low Energy (BLE)}
\term{BLE} tạo mạng cá nhân công suất thấp (PAN), phù hợp wearable, beacon, cảm biến ở gần điện thoại hoặc gateway. BLE phân biệt \emph{advertising/scanning} với \emph{connected mode}, và mô hình dịch vụ GATT/attribute. Chu kỳ quảng bá, connection interval, MTU, PHY, công suất phát và cơ chế sleep quyết định năng lượng. BLE không trực tiếp mang IP trong mô hình triển khai thường gặp; điện thoại hoặc gateway có thể chuyển dữ liệu lên backend. Bảng so sánh Bluetooth Classic và BLE ở slide 27 thuộc một số cấu hình thế hệ trước, không phải giới hạn chung của mọi phiên bản BLE hiện nay.\srcslide{26--27}
\begin{keybox}
Với cảm biến chạy pin, cần xét \emph{năng lượng trung bình} thay vì chỉ dòng phát cực đại: \[I_{\mathrm{avg}}\approx\sum_{k=1}^{n}I_k\frac{t_k}{T},\qquad T_{\mathrm{pin}}\approx\frac{C_{\mathrm{hữu\ dụng}}}{I_{\mathrm{avg}}}.\]
Các ước lượng trên bỏ qua nhiệt độ, tự xả, suy giảm pin và hiệu suất chuyển đổi nguồn; dùng để hiểu xu hướng, không phải tính tuổi thọ chính xác.
\end{keybox}
\subsection{Zigbee và mạng mesh IEEE 802.15.4}
\term{Zigbee} hỗ trợ các cấu trúc \emph{star}, \emph{tree} và \emph{mesh}, tiêu thụ thấp, có cơ chế bảo mật AES-128, tránh va chạm, ACK và thử truyền lại tùy tầng. Mạng có \emph{coordinator}, \emph{router} chuyển tiếp và \emph{end device}; thiết bị cuối có thể ngủ để tiết kiệm điện. Slide ghi khả năng khoảng 65.000 nút: đây là trần lý thuyết của một không gian địa chỉ, không đồng nghĩa một mạng thực tế có thể phục vụ số nút đó với độ trễ và thông lượng thỏa mãn.\srcslide{28--30}
\begin{figure}[htbp]\centering
\begin{tikzpicture}[every node/.style={circle,draw=navy,minimum size=6mm,inner sep=1pt},line width=.7pt]
\node[fill=accent!50] (c) at (0,0) {C};
\node[fill=blue!30] (r1) at (-2,1) {R};\node[fill=blue!30] (r2) at (2,1) {R};\node[fill=blue!30] (r3) at (0,2) {R};
\node[fill=pale] (e1) at (-3,2.1) {E};\node[fill=pale] (e2) at (-1.9,2.6) {E};\node[fill=pale] (e3) at (3,2.0) {E};\node[fill=pale] (e4) at (1,3) {E};
\draw (c)--(r1)--(r3)--(r2)--(c) (r1)--(r2) (r1)--(e1) (r1)--(e2) (r2)--(e3) (r3)--(e4);
\node[draw=none,rectangle,anchor=west] at (4.2,2.1) {C: Coordinator};
\node[draw=none,rectangle,anchor=west] at (4.2,1.5) {R: Router};
\node[draw=none,rectangle,anchor=west] at (4.2,.9) {E: End device};
\end{tikzpicture}\caption{Minh họa mesh Zigbee: router có thể chuyển tiếp qua nhiều chặng.}
\end{figure}
\subsection{Z-Wave}
\term{Z-Wave} chủ yếu dùng trong tự động hóa nhà ở: công tắc, khóa cửa, cảm biến an ninh, bộ điều khiển. Bài giảng nêu các dải xấp xỉ 868.42 MHz (châu Âu) và 908.42 MHz (Hoa Kỳ), một số mức dữ liệu 9.6/40/200 kbps và khoảng cách vài chục mét trong nhà. Băng tần cụ thể phụ thuộc thị trường, phiên bản và quy định quản lý phổ; không có một tần số hợp pháp chung toàn thế giới.\srcslide{31}
\subsection{LPWAN không dùng mạng di động: Weightless, Sigfox và LoRaWAN}
\term{Weightless} là họ tiêu chuẩn LPWAN gồm các biến thể W, P và N, khác nhau về phổ tần, hướng truyền và bài toán triển khai. Weightless-W nhắc đến TV white space; Weightless-P hướng tới hai chiều và mạng riêng. Khả năng thương mại/độ phổ biến thực tế phải được kiểm tra theo địa bàn.
\term{Sigfox} là một hệ sinh thái LPWAN phục vụ gói dữ liệu nhỏ, không thường xuyên và ưu tiên pin lâu. Kết nối phụ thuộc vùng có mạng và mô hình dịch vụ. Không thích hợp luồng video hoặc dữ liệu điều khiển thời gian thực dung lượng cao.
\term{LoRa} là kỹ thuật điều chế/vật lý vô tuyến tầm xa, còn \term{LoRaWAN} là giao thức mạng/MAC, quản lý thiết bị và hạ tầng gateway/network server trên LoRa. LoRaWAN hỗ trợ cấu hình lớp thiết bị A/B/C (đánh đổi năng lượng và độ trễ downlink), trao đổi hai chiều và cơ chế khóa bảo vệ mạng/ứng dụng. Băng thông thấp, giới hạn airtime và duty-cycle tùy khu vực là các ràng buộc thiết kế.\srcslide{32--34}
\begin{figure}[htbp]\centering
\begin{tikzpicture}[>=Latex,font=\footnotesize,node distance=12mm and 12mm]
\node[draw,rounded corners,fill=pale] (e1) {Node LoRa A};
\node[draw,rounded corners,fill=pale,below=9mm of e1] (e2) {Node LoRa B};
\node[draw,rounded corners,fill=blue!15,right=27mm of e1,yshift=-5mm] (gw) {Gateway};
\node[draw,rounded corners,fill=blue!15,right=of gw] (ns) {Network server};
\node[draw,rounded corners,fill=accent!15,right=of ns] (app) {Ứng dụng};
\draw[->,thick] (e1)--(gw);\draw[->,thick] (e2)--(gw);\draw[<->,thick] (gw)--node[above]{IP}(ns);\draw[<->,thick] (ns)--(app);
\end{tikzpicture}\caption{Phân biệt liên kết vô tuyến LoRa với kiến trúc mạng LoRaWAN.}
\end{figure}
\subsection{Li-Fi, mạng di động, 5G, LTE-M và NB-IoT}
\term{Li-Fi} truyền dữ liệu bằng ánh sáng (ví dụ LED), phù hợp một số không gian giới hạn. Vật chắn/ánh sáng môi trường và yêu cầu bố trí thu phát ảnh hưởng mạnh đến kết nối; không phải mọi hệ thống Li-Fi đều cần đường nhìn thẳng tuyệt đối, nhưng không truyền xuyên tường như sóng RF điển hình.
\term{Cellular} gồm các thế hệ 3G/4G/LTE/5G, dùng hạ tầng trạm gốc và lõi mạng nhà cung cấp, thuận tiện cho thiết bị di động, triển khai diện rộng. Thiết bị cần SIM/eSIM theo hệ thống mạng, thường có chi phí thuê bao và dòng đỉnh khi phát đáng chú ý. \term{5G} hỗ trợ nhiều kịch bản từ băng rộng di động đến IoT công nghiệp; các đặc trưng thực tế phụ thuộc băng tần, cấu hình mạng và vùng phủ.
\term{LTE-M (Cat-M1)} là công nghệ 3GPP LPWAN trên LTE, cân bằng di động, dữ liệu và pin, có thể thích hợp theo dõi tài sản. \term{NB-IoT} tối ưu cho dữ liệu nhỏ, phủ sâu và tiêu thụ thấp ở thiết bị ít dịch chuyển; độ trễ và mức hỗ trợ chuyển vùng/mobility tùy mạng. Phát biểu ``pin trên 10 năm'' chỉ đúng với chế độ ngủ, chu kỳ truyền, dung lượng pin và vùng phủ phù hợp, không phải tuổi thọ mặc định.\srcslide{35--40}
\begin{table}[htbp]\centering\small
\caption{Ma trận định tính để lựa chọn công nghệ (tổng hợp từ slide; mang tính định hướng)}
\begin{tabularx}{\textwidth}{@{}p{24mm}p{25mm}p{25mm}X@{}}\toprule
Công nghệ & Phạm vi điển hình & Năng lượng & Trường hợp sử dụng\\\midrule
NFC & Cực gần & Nhỏ cho thẻ thụ động & Chạm thẻ, mở khóa\\
BLE & Gần & Thấp nếu tối ưu sleep & Đồng hồ, wearable\\
Wi-Fi & LAN & Trung bình/cao & Camera, gateway, thiết bị có nguồn\\
Zigbee & PAN/mesh & Thấp & Cảm biến nhà thông minh\\
LoRaWAN & Xa & Thấp, tốc độ nhỏ & Đồng hồ, nông nghiệp\\
LTE-M & Diện rộng & Tùy mạng/duty cycle & Thiết bị di động có SIM\\
NB-IoT & Diện rộng & Thấp với chu kỳ thưa & Công tơ cố định, cảnh báo\\
UWB & Gần & Tùy chế độ & Đo khoảng cách, indoor tracking\\\bottomrule
\end{tabularx}
\end{table}
\section{Bảo mật và định tuyến trong IoT}
\subsection{An toàn thông tin và phân tầng bảo mật}
Trong IoT, ba thuộc tính cơ bản là \emph{confidentiality} (bí mật dữ liệu), \emph{integrity} (toàn vẹn), \emph{availability} (sẵn sàng). Cần thêm \emph{authentication} (xác thực chủ thể) và \emph{authorization} (phân quyền hành động). Mã hóa mà không kiểm tra tính xác thực và quản lý khóa đúng cách vẫn có thể bị giả mạo hoặc replay.\srcslide{42}
\begin{description}[style=nextline]
\item[IPsec.] Bộ cơ chế bảo vệ ở lớp IP, có thể cung cấp xác thực, bảo toàn vẹn và mã hóa theo cấu hình; sử dụng với IPv4 hoặc IPv6, không tự động bật chỉ vì dùng IPv6.
\item[TLS.] Bảo mật phiên kết nối thường đặt trên TCP, dùng cho HTTPS, MQTT over TLS và các dịch vụ ứng dụng khác. Xác thực chứng chỉ và khóa là điều kiện quan trọng.
\item[DTLS.] Tương tự TLS nhưng thiết kế cho datagram (thường UDP), thường được đề cập với CoAP. Ngoài ra, một số triển khai CoAP có thể dùng bảo mật đối tượng và vận chuyển bảo mật khác.
\item[IEEE 1888.3.] Slide giới thiệu chuẩn an toàn liên quan họ IEEE 1888 cho mạng điều khiển cộng đồng xanh, nhằm hạn chế truy cập không hợp lệ và bảo vệ dữ liệu.
\end{description}
\begin{figure}[htbp]\centering
\begin{tikzpicture}[font=\footnotesize,>=Latex]
\node[draw,fill=pale,minimum width=49mm,minimum height=10mm] (a) at (0,1.5) {Ứng dụng: MQTT, HTTP};
\node[draw,fill=pale,minimum width=49mm,minimum height=10mm] (b) at (0,.3) {TCP + TLS};
\node[draw,fill=pale,minimum width=49mm,minimum height=10mm] (c) at (0,-.9) {IP + Wi-Fi/cellular};
\node[draw,fill=accent!12,minimum width=49mm,minimum height=10mm] (d) at (6.3,1.5) {Ứng dụng: CoAP};
\node[draw,fill=accent!12,minimum width=49mm,minimum height=10mm] (e) at (6.3,.3) {UDP + DTLS};
\node[draw,fill=accent!12,minimum width=49mm,minimum height=10mm] (f) at (6.3,-.9) {IP + 802.15.4/khác};
\draw[->] (a)--(b)--(c);\draw[->] (d)--(e)--(f);
\end{tikzpicture}\caption{Ví dụ các ngăn xếp giao thức được bảo vệ.}
\end{figure}
\subsection{Định tuyến RPL (Routing Protocol for Low-Power and Lossy Networks)}
\term{Định tuyến} chọn đường đưa gói tin từ nguồn đến đích qua một hay nhiều hop. IoT thường có nút chạy pin, liên kết không ổn định, tốc độ thấp và số nút lớn. \term{RPL} là giao thức định tuyến IPv6 hướng tới mạng năng lượng thấp và suy hao cao (LLN), phù hợp một số mạng 802.15.4. RPL tổ chức một \emph{DODAG} (Destination-Oriented Directed Acyclic Graph) hướng về root; mỗi nút tính \emph{rank} theo hàm mục tiêu, có thể dựa trên số hop, chỉ số liên kết ETX hoặc thước đo khác.\srcslide{43--45}
\begin{figure}[htbp]\centering
\begin{tikzpicture}[>=Latex,node distance=13mm,font=\footnotesize,every node/.style={draw,circle,fill=pale,minimum size=9mm}]
\node[fill=accent!35] (root) at (0,2.7) {R};
\node (a) at (-2,1.2) {A};\node (b) at (2,1.2) {B};
\node (c) at (-3,0) {C};\node (d) at (-.8,0) {D};\node (e) at (2.6,0) {E};
\draw[->,thick,blue] (c)--(a)--(root);\draw[->,thick,blue] (d)--(a);\draw[->,thick,blue] (e)--(b)--(root);
\draw[densely dashed] (d)--(b);
\node[draw=none,rectangle] at (4.0,2.7) {R = Root/Gateway};
\node[draw=none,rectangle] at (4.0,2.1) {Mũi tên: hướng uplink};
\end{tikzpicture}\caption{DODAG minh họa trong RPL; liên kết nét đứt là tuyến thay thế có thể cân nhắc.}
\end{figure}
\begin{keybox}Không đồng nhất \textbf{routing protocol} với \textbf{network topology}: mesh mô tả các liên kết có thể có, còn RPL quyết định con đường chuyển gói trên một mạng tương ứng.\end{keybox}
\section{Bốn loại mạng IoT (Types of IoT Network)}
Bài giảng phân nhóm phổ biến theo vùng phủ và tổ chức kết nối thành \term{Cellular}, \term{LAN/PAN}, \term{LPWAN} và \term{Mesh}. Đây là một cách phân loại thực dụng, \emph{không loại trừ lẫn nhau}: mesh là topology, LPWAN là lớp đặc trưng về năng lượng và diện phủ; mạng có thể đồng thời là mạng mesh thuộc PAN.\srcslide{46--54}
\subsection{Cellular networks}
Dùng hạ tầng di động và các trạm gốc của nhà mạng. Ưu điểm: vùng phủ diện rộng, hỗ trợ thiết bị di chuyển, quản lý thuê bao/hạ tầng sẵn có. Nhược điểm: có thể mất sóng tại tầng hầm, hầm mỏ, tủ điện kín; chi phí kết nối; đỉnh dòng phát và nhu cầu quản lý năng lượng. LTE-M và NB-IoT là hai lựa chọn Cellular IoT được nêu trong bài giảng.\srcslide{47--49}
\subsection{Local Area Network và Personal Area Network}
\term{LAN} kết nối các thiết bị trong phạm vi địa điểm như phòng học, văn phòng, xưởng hoặc hộ gia đình; ví dụ Ethernet/Wi-Fi. \term{PAN} nhấn mạnh phạm vi quanh cá nhân; BLE là ví dụ thường gặp. LAN/PAN có ưu điểm dễ triển khai và chi phí tương đối thấp nếu đã có hạ tầng; nhưng vùng phủ không rộng và chịu tác động nhiễu, vật chắn, tải mạng. Không thể mặc định tất cả LAN/PAN ``không tin cậy'': chất lượng tùy chuẩn và thiết kế hệ thống.\srcslide{50--51}
\subsection{LPWAN}
Mạng \term{Low-Power Wide-Area Network} nhắm tới cảm biến công suất thấp trao đổi payload nhỏ trên khoảng cách lớn; đánh đổi bằng tốc độ hạn chế, độ trễ có thể cao và giới hạn số lần truyền. LoRaWAN, Sigfox, LTE-M và NB-IoT là những công nghệ có thể xếp vào nhóm LPWAN, dù khác nhau về hạ tầng và giao thức.\srcslide{52}
\subsection{Mesh networks}
Trong mesh, các router/node có thể chuyển tiếp gói qua nhiều chặng, mở rộng vùng phủ ngoài liên kết trực tiếp tới gateway. Ưu điểm: đường đi linh hoạt và có thể phục hồi khi một tuyến yếu. Nhược điểm: số hop làm tăng độ trễ, overhead định tuyến và số lần phát vô tuyến; router trung gian chạy pin có thể nhanh hết năng lượng. Không phải mọi thiết bị đầu cuối đều làm nhiệm vụ chuyển tiếp.\srcslide{53--54}
\begin{table}[htbp]\centering\small\caption{So sánh khái quát bốn nhóm mạng trong slide}
\begin{tabularx}{\textwidth}{@{}p{23mm}X X@{}}\toprule
Nhóm & Ưu điểm tiêu biểu & Hạn chế cần xét\\\midrule
Cellular & Rộng, có thể hỗ trợ di động & Phụ thuộc phủ sóng, SIM, chi phí và dòng đỉnh\\
LAN/PAN & Hạ tầng phổ biến, dễ tích hợp & Vùng phủ và nhiễu nội bộ\\
LPWAN & Xa, pin lâu nếu gửi thưa & Payload và thông lượng hạn chế\\
Mesh & Nhiều đường, mở rộng bằng hop & Routing overhead, độ trễ và router trung gian\\\bottomrule
\end{tabularx}
\end{table}
\section{Giao thức trao đổi dữ liệu IoT (IoT Data Protocols)}
\subsection{Hai mô hình trao đổi chính}
\term{Request/response} (yêu cầu/đáp ứng): client gửi yêu cầu cụ thể đến server, server trả kết quả; thường gặp với HTTP, CoAP. \term{Publish/subscribe} (xuất bản/đăng ký): publisher gửi bản tin vào một \emph{topic}, broker chuyển cho subscriber đã đăng ký topic. Mô hình thứ hai giảm sự ràng buộc trực tiếp giữa bên phát và bên nhận, phù hợp nhiều nguồn và nhiều đối tượng nghe dữ liệu. \term{Subscribe} ở MQTT là đăng ký chủ đề nhận bản tin, \emph{không đồng nghĩa} thủ tục đăng ký định danh thiết bị vào mạng hoặc xác thực tài khoản như slide 56 có thể gây hiểu nhầm.\srcslide{55--59}
\begin{figure}[htbp]\centering
\begin{tikzpicture}[>=Latex,font=\footnotesize]
\node[draw,rounded corners,fill=pale] (cl) at (0,2.3) {Client};\node[draw,rounded corners,fill=pale] (sv) at (4.4,2.3) {Server};
\draw[->,blue,thick] (cl.north east) -- node[above]{Request} (sv.north west);\draw[->,accent,thick] (sv.south west) -- node[below]{Response} (cl.south east);
\node[draw,rounded corners,fill=pale] (pub) at (0,-.5) {Publisher};\node[draw,rounded corners,fill=blue!14] (br) at (4.4,-.5) {Broker};\node[draw,rounded corners,fill=pale] (sub) at (8.5,-.5) {Subscriber};
\draw[->,blue,thick] (pub)--node[above]{PUBLISH}(br);\draw[->,blue,thick] (br)--node[above]{Topic match}(sub);
\draw[->,accent,thick] (sub.south) .. controls (8,-1.9) and (5,-1.9) .. node[below]{SUBSCRIBE topic} (br.south);
\end{tikzpicture}\caption{Request/response (trên) và publish/subscribe qua broker (dưới).}
\end{figure}
\subsection{HTTP: Hypertext Transfer Protocol}
HTTP là giao thức ứng dụng phổ biến cho web và REST API; giao tiếp theo yêu cầu/đáp ứng, có khả năng thương lượng định dạng nội dung và về bản chất \term{stateless} (mỗi yêu cầu độc lập ở cấp giao thức). Các phương thức điển hình gồm GET (đọc), POST (gửi/tạo hoặc xử lý), PUT (tạo/thay thế), PATCH (cập nhật từng phần), DELETE (xóa). Cần phân biệt tính \emph{idempotent}: GET, PUT, DELETE được định nghĩa có tính chất này, còn POST thông thường không. HTTP có thể dùng kết nối tồn tại lâu, streaming, HTTP/2 hoặc HTTP/3 tùy hạ tầng; vì vậy phát biểu ở slide rằng HTTP luôn chỉ là giao tiếp ``một chiều'' và không thể mở rộng là cách đơn giản hóa.\srcslide{60--64}
\begin{verbatim}
GET /api/sensors/thermometer-01 HTTP/1.1
Host: example.com
Accept: application/json

HTTP/1.1 200 OK
Content-Type: application/json

{"id":"thermometer-01","temperature_c":28.4}
\end{verbatim}
\term{Khi dùng HTTP:} dashboard, cấu hình thiết bị, API quản trị và tích hợp dịch vụ. Thiết bị chạy pin muốn truyền đều đặn mỗi vài giây qua kết nối di động có thể phải trả overhead lớn hơn so với MQTT/CoAP trong một số trường hợp. TLS tạo HTTPS để bảo mật dữ liệu trong vận chuyển.
\subsection{TCP so với UDP}
\term{TCP} hướng kết nối, truyền chuỗi byte theo thứ tự, phát hiện mất/nhận lại gói và điều khiển luồng/tắc nghẽn. \term{UDP} không thiết lập kết nối tại tầng vận chuyển và không bảo đảm giao nhận/thứ tự; ứng dụng có thể tự bổ sung ACK, timeout, retransmission. UDP có header cơ bản 8 byte; TCP có header tối thiểu 20 byte (không kể IP, tùy chọn và mã hóa). Không nên khẳng định ``UDP luôn nhanh hơn TCP'' trong mọi bài toán: độ tin cậy và chi phí toàn hệ thống quyết định hiệu quả.\srcslide{64--65}
\begin{table}[htbp]\centering\small\caption{TCP và UDP đối với giao thức ứng dụng IoT}
\begin{tabularx}{\textwidth}{@{}p{32mm}X X@{}}\toprule
Thuộc tính & TCP & UDP\\\midrule
Kết nối & Có & Không\\
Đảm bảo truyền tải & Có truyền lại và thứ tự byte & Không ở tầng UDP\\
Header tối thiểu & 20 byte & 8 byte\\
Ví dụ & HTTP/1.1, MQTT, WebSocket & CoAP mặc định\\
Bảo mật thường gặp & TLS & DTLS\\\bottomrule
\end{tabularx}
\end{table}
\subsection{CoAP: Constrained Application Protocol}
CoAP được thiết kế cho thiết bị và mạng hạn chế tài nguyên; theo mô hình REST gần HTTP với \texttt{GET}, \texttt{POST}, \texttt{PUT}, \texttt{DELETE}, URI, content format, proxy và cache. CoAP thường chạy trên UDP; phiên bản tiêu chuẩn còn có khả năng hoạt động qua TCP và WebSocket trong các quy định mở rộng. Trong trường hợp CoAP/UDP, bản tin có bốn loại: \term{CON} (confirmable, yêu cầu xác nhận), \term{NON} (non-confirmable), \term{ACK} (xác nhận) và \term{RST} (reset). Định danh bản tin 16-bit giúp nhận diện trùng lặp trong ngữ cảnh áp dụng; token gắn phản hồi với yêu cầu.\srcslide{65--70}
\begin{figure}[htbp]\centering
\begin{tikzpicture}[>=Latex,font=\footnotesize]
\node (a) at (0,0) {Client};\node (b) at (8,0) {Server};
\draw[dashed] (0,-.3)--(0,-3);\draw[dashed] (8,-.3)--(8,-3);
\draw[->,thick,blue] (0,-.8)--node[above]{CON GET, MID=100, Token=AA}(8,-1.25);
\draw[<-,thick,accent] (0,-2.2)--node[above]{ACK 2.05 Content, MID=100}(8,-1.75);
\end{tikzpicture}\caption{Ví dụ bản tin CoAP confirmable với phản hồi gộp trong ACK.}
\end{figure}
\begin{keybox}CoAP truyền trên UDP không có bảo đảm mặc định ở tầng vận chuyển, nhưng \textbf{CON + ACK + cơ chế retransmission ở tầng CoAP} cung cấp độ tin cậy ở cấp trao đổi thông điệp. Không nên kết luận ``CoAP không có ACK'' chỉ từ việc dùng UDP.\end{keybox}
\subsection{MQTT: Message Queuing Telemetry Transport}
MQTT là giao thức nhắn tin nhẹ theo publish/subscribe, ứng dụng nhiều với telemetry và M2M. Các thành phần: \term{publisher}, \term{subscriber}, \term{broker}, \term{topic}, \term{payload}. Client thường tạo phiên TCP với broker; kết nối có thể được bảo vệ bằng TLS. Ví dụ cảm biến xuất bản lên \texttt{factory/line1/temperature}, dashboard đăng ký \texttt{factory/+/temperature}; ký tự \texttt{+} là wildcard một tầng, \texttt{\#} là wildcard nhiều tầng ở vị trí hợp lệ trong topic filter.\srcslide{71--75}
\begin{figure}[htbp]\centering
\begin{tikzpicture}[font=\footnotesize,>=Latex,node distance=21mm]
\node[draw,rounded corners,fill=pale] (sensor) {Cảm biến};
\node[draw,rounded corners,fill=blue!20,right=of sensor] (broker) {MQTT Broker};
\node[draw,rounded corners,fill=pale,above right=9mm and 23mm of broker] (d1) {Dashboard};
\node[draw,rounded corners,fill=pale,below right=9mm and 23mm of broker] (d2) {Dịch vụ lưu CSDL};
\draw[->,thick,blue] (sensor)--node[above]{PUBLISH}(broker);
\draw[->,thick,blue] (broker)--node[above,sloped]{PUBLISH}(d1);
\draw[->,thick,blue] (broker)--node[below,sloped]{PUBLISH}(d2);
\draw[->,dashed] (d1.west) to[bend right=35] node[below]{SUBSCRIBE} (broker.north);
\end{tikzpicture}\caption{MQTT broker phân phối dữ liệu theo topic đến nhiều subscriber.}
\end{figure}
\subsubsection{Bốn giai đoạn và thông điệp}
Slide tóm tắt một phiên giao tiếp thành \emph{connection}, \emph{authentication}, \emph{communication}, \emph{termination}. Ở cấp giao thức, có các gói như CONNECT/CONNACK, PUBLISH, SUBSCRIBE/SUBACK, PINGREQ/PINGRESP, DISCONNECT. Tên giai đoạn ``authentication'' trong slide là mô tả logic: xác thực có thể là một phần của quá trình CONNECT, có cơ chế nâng cao tùy phiên bản và cấu hình broker.
\subsubsection{Ba cấp QoS}
\begin{itemize}
\item \term{QoS 0 — At most once:} gửi theo mức nỗ lực tối đa, không có retransmission của MQTT cho gói publish này. Có thể mất tin.
\item \term{QoS 1 — At least once:} dùng PUBACK; có thể phát lại, subscriber có khả năng nhận trùng. Ứng dụng nên thiết kế idempotent nếu cần.
\item \term{QoS 2 — Exactly once:} chuỗi xác nhận nhiều bước (PUBREC/PUBREL/PUBCOMP) nhằm tránh phát MQTT PUBLISH trùng trong ngữ cảnh trao đổi giữa các bên MQTT; tăng overhead. Không đồng nghĩa bảo đảm ``mỗi tác động vật lý diễn ra đúng một lần'' ở toàn hệ thống khi mất nguồn hoặc có lỗi nghiệp vụ.
\end{itemize}
\begin{table}[htbp]\centering\small\caption{So sánh MQTT QoS}
\begin{tabularx}{\textwidth}{@{}p{16mm}p{35mm}X@{}}\toprule
QoS & Cam kết phân phối & Đánh đổi / ví dụ\\\midrule
0 & At most once & Ít overhead; telemetry liên tục có thể chấp nhận mất một vài mẫu\\
1 & At least once & Có thể trùng; báo trạng thái quan trọng, áp dụng ID để khử trùng\\
2 & Exactly once ở cấp MQTT & Overhead cao; chỉ chọn khi thật cần cơ chế phân phối chặt chẽ\\\bottomrule
\end{tabularx}
\end{table}
\begin{note}
Các thuộc tính \emph{retained message}, \emph{persistent session/session expiry}, \emph{last will} và \emph{keep alive} (mở rộng kỹ thuật) đặc biệt hữu ích khi thiết bị IoT mất kết nối. QoS không thay thế cho sao lưu ở SD, hàng đợi bền vững và xác nhận nghiệp vụ của ứng dụng.
\end{note}
\subsection{WebSocket}
WebSocket mở kết nối hai chiều liên tục (full-duplex), thường bắt đầu bằng HTTP Upgrade trên HTTP/1.1, dùng TCP và có thể được mã hóa bằng TLS (\texttt{wss://}). Phù hợp dashboard thời gian thực, điều khiển trình duyệt và ứng dụng web tương tác. Không phải bản thân WebSocket được tối ưu năng lượng cho MCU chạy pin; cấu hình heartbeat, duy trì kết nối và dữ liệu gửi ảnh hưởng đáng kể. WebSocket khác HTTP request/response nhưng không phải ``HTTP nhanh hơn'' hay ``WebSocket luôn nặng'' trong mọi ngữ cảnh.\srcslide{76--77}
\subsection{AMQP và XMPP}
\term{AMQP (Advanced Message Queuing Protocol)} là giao thức nhắn tin hướng trung gian có cơ chế quản lý hàng đợi và định tuyến thông điệp. Trong mô hình AMQP 0-9-1 phổ biến, \emph{exchange} nhận và định tuyến đến \emph{queue} theo binding/routing key; \emph{consumer} lấy thông điệp từ queue. Chi tiết ``exactly once'' phụ thuộc đặc tả, broker và cơ chế xử lý ở ứng dụng, không nên coi là bảo đảm mặc định của AMQP. Bảo mật thường triển khai TLS và cơ chế xác thực phù hợp.\srcslide{78}
\term{XMPP (Extensible Messaging and Presence Protocol)} là giao thức liên lạc dùng XML, hỗ trợ định danh, phiên và trao đổi thông điệp qua máy chủ; dùng TCP theo triển khai phổ biến. Khả năng pub/sub được bổ sung qua các XMPP Extension Protocols, không phải mọi triển khai đều tự có. XMPP phù hợp các tình huống cần hiện diện và nhắn tin phân tán, nhưng XML có overhead đối với nút tài nguyên hạn chế.\srcslide{79}
\begin{table}[htbp]\centering\scriptsize\caption{Tổng hợp giao thức ứng dụng được giảng dạy}
\begin{tabularx}{\textwidth}{@{}p{19mm}p{26mm}p{22mm}X@{}}\toprule
Giao thức & Mô hình & Nền vận chuyển phổ biến & Điểm nổi bật\\\midrule
HTTP & Request/response & TCP/TLS; HTTP/3 qua QUIC & REST, web API, quản trị\\
CoAP & REST request/response & UDP/DTLS & Overhead thấp, CON/ACK\\
MQTT & Publish/subscribe & TCP/TLS & Broker, topic, QoS 0/1/2\\
WebSocket & Hai chiều liên tục & TCP/TLS & Realtime browser/server\\
AMQP & Message queue/pub-sub & TCP/TLS & Queue, exchange, routing\\
XMPP & Messaging, các mở rộng pub-sub & TCP/TLS & XML, presence, federation\\\bottomrule
\end{tabularx}
\end{table}
\section{Áp dụng vào thiết kế mạng IoT}
\subsection{Bài toán 1: cảm biến nhà thông minh}
Hệ thống gồm hàng chục cảm biến cửa/nhiệt độ, đặt trong một căn nhà, chạy pin, gửi vài byte mỗi phút. Lựa chọn có cơ sở là Zigbee mesh hoặc BLE đến gateway, tùy thiết bị và yêu cầu; phần gateway dùng Ethernet/Wi-Fi đến MQTT broker hoặc dịch vụ quản lý. Không chọn Wi-Fi trực tiếp theo mặc định nếu thời lượng pin là ưu tiên; không cần dùng 5G hay UWB nếu không có yêu cầu di động/định vị.
\subsection{Bài toán 2: giám sát nông nghiệp diện rộng}
Nếu có hạ tầng LoRaWAN hợp lệ tại khu vực, các nút đo ẩm đất gửi vài bản tin mỗi giờ tới gateway là một phương án. Nếu thiếu gateway nhưng có phủ sóng NB-IoT, có thể dùng modem tế bào kèm SIM. Cần xem xét dung lượng payload, chu kỳ nhận downlink, nguồn cấp và chi phí dài hạn. Đối với dữ liệu cần bảo đảm, lưu đệm mẫu cục bộ và đánh số thứ tự.
\subsection{Bài toán 3: thiết bị đeo giám sát và báo nguy}
Thiết bị đeo thu IMU, truyền trạng thái quan trọng theo chu kỳ đến gateway ở gần; lớp giao tiếp có thể là BLE hoặc RF phù hợp và lớp ứng dụng sử dụng MQTT qua gateway. Mất kết nối cần có cơ chế lưu bộ nhớ, gửi lại, thông báo trạng thái liên kết. QoS 1 cần idempotency để không xử lý một cảnh báo lặp thành nhiều sự kiện. Cảnh báo an toàn tính mạng cần thiết kế dự phòng và kiểm chứng thực nghiệm; một giao thức không thể tự đảm bảo an toàn hệ thống.
\section{Công thức và thuật ngữ trọng tâm}
\begin{align}
N_{\mathrm{IPv4}}&=2^{32}\approx4.29\times 10^{9}\\
N_{\mathrm{IPv6}}&=2^{128}\approx3.40\times10^{38}\\
R_{\mathrm{hiệu\ dụng}}&=\frac{\mathrm{Payload\ bits\ nhận\ đúng}}{\mathrm{Thời\ gian\ quan\ sát}}\\
\mathrm{PDR}&=\frac{\mathrm{Số\ gói\ hợp\ lệ\ đến\ đích}}{\mathrm{Số\ gói\ phát}}\\
I_{\mathrm{avg}}&=\frac{1}{T}\int_0^T I(t)\,dt.
\end{align}
\noindent Các đại lượng nên được đo trong điều kiện thử nghiệm xác định; tốc độ dữ liệu danh nghĩa (PHY bit rate) không phải dữ liệu hữu ích tại ứng dụng.
\begin{longtable}{@{}p{32mm}p{103mm}@{}}
\caption{Bảng thuật ngữ viết tắt cần nhớ}\\\toprule
Thuật ngữ & Ý nghĩa\\\midrule\endfirsthead
\toprule Thuật ngữ & Ý nghĩa\\\midrule\endhead
IoT & Internet of Things — Internet vạn vật\\
M2M/H2M/M2H/H2H & Máy--máy / người--máy / máy--người / người--người\\
PAN/LAN & Personal/Local Area Network\\
LPWAN & Low-Power Wide-Area Network\\
IPv4/IPv6 & Internet Protocol phiên bản 4 / 6\\
NFC/UWB & Near-Field Communication / Ultra-Wideband\\
BLE & Bluetooth Low Energy\\
PLC & Power-Line Communication\\
RPL & Routing Protocol for Low-Power and Lossy Networks\\
DODAG & Destination-Oriented Directed Acyclic Graph\\
TLS/DTLS & Transport Layer Security / Datagram TLS\\
CoAP & Constrained Application Protocol\\
MQTT & Message Queuing Telemetry Transport\\
QoS & Quality of Service\\
AMQP & Advanced Message Queuing Protocol\\
XMPP & Extensible Messaging and Presence Protocol\\
CON/NON/ACK/RST & Các loại bản tin CoAP: confirmable / nonconfirmable / acknowledgment / reset\\
\bottomrule
\end{longtable}
\section{Câu hỏi tự kiểm tra và đáp án ngắn}
\begin{enumerate}[leftmargin=7mm,itemsep=7pt]
\item \textbf{Hỏi:} M2M khác M2H ở điểm nào? \textbf{Đáp:} M2M truyền thông tin giữa các máy; M2H là máy tạo thông báo cho người.
\item \textbf{Hỏi:} Chênh lệch độ dài địa chỉ IPv4 và IPv6? \textbf{Đáp:} IPv4 32 bit, IPv6 128 bit, tức IPv6 dài gấp 4 lần về số bit.
\item \textbf{Hỏi:} Vì sao địa chỉ IP không thay thế xác thực? \textbf{Đáp:} IP là địa chỉ định tuyến, không chứng minh danh tính hoặc quyền truy cập.
\item \textbf{Hỏi:} LoRa khác LoRaWAN thế nào? \textbf{Đáp:} LoRa là lớp vật lý điều chế; LoRaWAN quy định các chức năng mạng/MAC và kiến trúc phía trên.
\item \textbf{Hỏi:} Khi nào nên cân nhắc BLE? \textbf{Đáp:} Truyền dữ liệu ít ở phạm vi gần, yêu cầu pin lâu và có gateway/điện thoại.
\item \textbf{Hỏi:} Zigbee có thể tạo các topology nào? \textbf{Đáp:} Star, tree, mesh (và các kiểu điểm-điểm ở tầng liên kết theo ngữ cảnh).
\item \textbf{Hỏi:} Vì sao mesh có thể tăng phủ sóng nhưng cũng tăng tiêu thụ? \textbf{Đáp:} Gói tin đi qua router trung gian, tạo thêm lần phát và overhead định tuyến.
\item \textbf{Hỏi:} RPL giải quyết vấn đề gì? \textbf{Đáp:} Định tuyến IPv6 trong mạng ít năng lượng và liên kết tổn hao.
\item \textbf{Hỏi:} HTTP và MQTT khác mô hình trao đổi ra sao? \textbf{Đáp:} HTTP thông thường request/response; MQTT publish/subscribe qua broker.
\item \textbf{Hỏi:} CoAP truyền trên UDP có thể bảo đảm ACK không? \textbf{Đáp:} Có, ở tầng CoAP bằng CON/ACK và cơ chế truyền lại.
\item \textbf{Hỏi:} QoS 1 có thể sinh bản tin trùng? \textbf{Đáp:} Có; ít nhất một lần, vì vậy cần ID và khử trùng khi xử lý.
\item \textbf{Hỏi:} QoS 2 có bảo đảm van nước chỉ mở đúng một lần trong mọi lỗi không? \textbf{Đáp:} Không; QoS 2 chỉ giải quyết ngữ nghĩa phân phối MQTT, nghiệp vụ cần bảo vệ riêng.
\item \textbf{Hỏi:} TLS và DTLS thường dùng với loại transport nào? \textbf{Đáp:} TLS trên giao thức dạng luồng (thường TCP), DTLS trên datagram (thường UDP).
\item \textbf{Hỏi:} Khi nào dùng WebSocket? \textbf{Đáp:} Dữ liệu hai chiều liên tục giữa browser/app với server, ví dụ dashboard thời gian thực.
\item \textbf{Hỏi:} Chọn NB-IoT hay Wi-Fi cho công tơ nơi xa? \textbf{Đáp:} Nếu có NB-IoT với vùng phủ và ngân sách phù hợp, NB-IoT thuận tiện hơn khi không có hạ tầng Wi-Fi; luôn khảo sát thực tế.
\end{enumerate}
\section{Bài tập vận dụng}
\textbf{Bài 1.} Một nút có chu kỳ 60 giây: ngủ 59.5 giây ở 0.02 mA, phát 0.5 giây ở 60 mA. Tính dòng trung bình và tuổi thọ lý thuyết khi dung lượng khả dụng là 2000 mAh.\par
\textbf{Hướng dẫn:} \(I_{avg}=0.02\cdot59.5/60+60\cdot0.5/60\approx0.5198\,\mathrm{mA}\). Thời gian \(2000/0.5198\approx3848\) giờ, khoảng 160 ngày (không kể hiệu suất nguồn, tự xả, nhiệt độ và lão hóa).\par\medskip
\textbf{Bài 2.} Thiết kế mạng cảm biến gồm 100 nút trong một nhà máy ba tầng, báo dữ liệu 30 giây/lần, báo động phải nhận dưới 3 giây; giải thích vì sao không thể chỉ dựa vào \emph{range} để chọn công nghệ.\par
\textbf{Gợi ý:} Cần tính dung lượng mạng, nhiều hop, nhiễu RF, retry, số router luôn bật, phân bổ gateway, kiểm tra độ trễ tail-percentile và bảo vệ chống mất kết nối.\par\medskip
\textbf{Bài 3.} Cảm biến nhiệt độ xuất MQTT QoS 1 lên \texttt{plant/a/temperature}. Một ứng dụng nhận trùng sau khi reconnect. Đề xuất định dạng payload và logic server.\par
\textbf{Gợi ý:} Gắn \texttt{device\_id}, \texttt{seq}, \texttt{ts}, \texttt{value}; xử lý khóa \texttt{device\_id+seq} duy nhất trong cơ sở dữ liệu và ghi nhận last sequence.
\section{Nguồn tham khảo và phạm vi tài liệu}
\begin{enumerate}[leftmargin=7mm]
\item Nguyễn Đắc Cử, \emph{Cơ sở và thiết kế hệ thống IoT — Chapter 3: Networking technologies in IoT}, slide bài giảng Đại học Phenikaa, 81 trang. \textbf{Nguồn chính} quyết định phạm vi và thứ tự chủ đề.
\item Qusay F. Hassan (ed.), \emph{Internet of Things A to Z: Technologies and Applications}, Wiley/IEEE Press, 2018, đặc biệt phần công nghệ lớp mạng, bảo mật và trao đổi dữ liệu.
\item Adrian McEwen và Hakim Cassimally, \emph{Designing the Internet of Things}, Wiley, 2014, chương Internet Principles và Prototyping Online Components.
\item Alice James, Avishkar Seth và Subhas Chandra Mukhopadhyay, \emph{IoT System Design: Project Based Approach}, Springer, ISBN 978-3-030-85863-6, 2022, các chương Wi-Fi, Bluetooth, LoRa.
\end{enumerate}
\begin{note}
Tài liệu là \textbf{bản tổng hợp học tập độc lập}, không phải bản chép nguyên văn slide hay sách. Các thông số trong slide là tài liệu môn học ở thời điểm biên soạn; khi làm dự án thực tế cần kiểm tra thông số thiết bị, đặc tả giao thức và quy định tần số tại nơi triển khai.
\end{note}